\documentclass[a4paper, 12pt]{article}
%\usepackage[polish]{babel}
\usepackage[cp1250]{inputenc}
%\usepackage[latin2]{inputenc}
%\usepackage{polski}
\usepackage{amsfonts}
\usepackage[T1]{fontenc}
\usepackage[dvips]{graphicx}

%\usepackage{pldoc}







\pagestyle{headings}
\textwidth      15.5cm
\oddsidemargin    .1cm
\evensidemargin   .1cm

\begin{document}

\thispagestyle{empty}

\begin{minipage}{5cm}
%\includegraphics[scale=0.3]{logo.eps}
\end{minipage}
\begin{minipage}{10cm}
\begin{center}
{\Large Salomon\\
System przetwarzania wiedzy\\}
\end{center}
\end{minipage}

\vspace*{0.5cm}

\hrulefill

\vspace*{1cm}


\begin{flushright}
{\Large \emph{Test Plan, ver 1.0} }
\end{flushright}

\vspace{7cm}

\begin{flushleft}
{ \Large Historia wersji }\[\]
\begin{tabular}{|p{2.5cm}||p{1.5cm}|p{6.5cm}|p{3.5cm}|}
  \hline
  Data & Wersja & Opis & Autorzy \\
  \hline \hline
  05-05-2005 & 1.0 & Pierwsza przymiarka do test�w & Tomasz W�sala\\
  \hline
\end{tabular}
\end{flushleft}



\[\]
\newpage
    \section{Wst�p}
        \subsection{Cel dokumentu}
            Dokument ma na celu przedstawienie podej�cia do test�w implementacji drzew decyzyjnych an platformie Salomon. Dotyczy on wersji pre-release.
    \section{Og�lny opis}
        Testy b�d� przygotowywane dla pierwszej wersji naszego projektu. B�dziemy korzystali z nast�puj�cych technologii:
        
        \begin{itemize}
            \item jUnit
            \item ewentualnie Gringer
        \end{itemize}
        Niekt�re testy, np. testy dzia�ania pluginu obs�uguj�cego wizualizacje, b�d� te� przeprowadzane r�cznie.

 \section{Plan test�w}
    Testy b�d� obejmowa�y sprawdzenie poprawno�ci:
    \begin{itemize}
        \item zrealizowanego algorytmu budowania drzew decyzyjnych.
        \item dzia�ania poszczeg�lnych klas, interfejs�w w tym b�d� zawarte interfejsy szkieletowe (np. rezlizuj�ce operacje dzia�ania na bazie danych).
        \item dzia�ania pluginu realizuj�cego budow� drzew decyzyjnych.
        \item dzia�ania atrapy pluginu tworz�cego wizualizacje.
        \item dzia�ania ca�o�ci zrealizowanej funkcjonalno�ci. 
    \end{itemize}

\end{document}
